\documentclass[aspectratio=169]{beamer}

% Theme selection
\usetheme{Madrid}
\usecolortheme{whale}

% Packages
\usepackage{graphicx}
\usepackage{tikz}
\usetikzlibrary{shapes,arrows,positioning}
\usepackage{booktabs}
\usepackage{listings}
\usepackage{xcolor}
\usepackage[fontset=fandol]{ctex}

% Custom colors
\definecolor{codegreen}{rgb}{0,0.6,0}
\definecolor{codegray}{rgb}{0.5,0.5,0.5}
\definecolor{codepurple}{rgb}{0.58,0.82}
\definecolor{backcolour}{rgb}{0.95,0.95,0.92}

% Code listing style
\lstdefinestyle{mystyle}{
    backgroundcolor=\color{backcolour},
    commentstyle=\color{codegreen},
    keywordstyle=\color{magenta},
    numberstyle=\tiny\color{codegray},
    stringstyle=\color{codepurple},
    basicstyle=\ttfamily\footnotesize,
    breakatwhitespace=false,
    breaklines=true,
    captionpos=b,
    keepspaces=true,
    numbers=left,
    numbersep=5pt,
    showspaces=false,
    showstringspaces=false,
    showtabs=false,
    tabsize=2
}
\lstset{style=mystyle}

\begin{document}

% Title Info
\title[塔科夫改枪优化器]{塔科夫改枪优化器}
\subtitle{基于约束规划的自动化配置方案优化}
\author{Imposs1ble嗨}
\date{\today}

\begin{frame}
    \titlepage
\end{frame}

% Section 1: Introduction
\section{背景介绍}

% Slide 3: The Problem
\begin{frame}{问题背景：复杂的武器改装系统}
    《逃离塔科夫》拥有业界最复杂的武器改装系统之一。

    \vspace{0.3cm}
    \textbf{玩家面临的主要挑战：}
    \begin{itemize}
        \item \textbf{配件数量庞大：}超过2500种配件，涉及后坐力、人机工效、重量、精准度等多项属性
        \item \textbf{嵌套兼容结构：}配件可安装于其他配件之上（导轨 $\to$ 转接座 $\to$ 瞄具）
        \item \textbf{组合空间爆炸：}一把AR-15有50余个槽位，有效配置方案多达数万种
        \item \textbf{性价比难以评估：}高价配件是否物有所值？
        \item \textbf{购买限制：}价格实时波动，且受商人等级和跳蚤市场解锁条件限制
    \end{itemize}

    \vspace{0.3cm}
    \textbf{现状：}玩家需要相当的时间来熟悉各个配件的属性，或直接照搬网上的"版本答案"，而不理解其背后的权衡逻辑或者盲目购买高价配件。
\end{frame}

% Slide 4: The Solution
\begin{frame}{解决方案：自动化配置优化}
    \textbf{项目目标：}
    在给定约束条件下，自动计算出数学意义上的最优改枪方案。

    \vspace{0.4cm}
    \begin{columns}
        \column{0.5\textwidth}
        \textbf{典型应用场景：}
        \begin{itemize}
            \item \textit{"预算20万卢布，如何改出最好的M4？"}
            \item \textit{"不限预算，AK的后坐力最低能达到多少？"}
            \item \textit{"当前商人等级下，能改出什么水平的武器？"}
            \item \textit{"增加10万预算，后坐力提升有多大？"}
            \item \textit{"不同预算档位下的最优改枪方案分别是什么？"}
        \end{itemize}

        \column{0.5\textwidth}
        \textbf{核心功能：}
        \begin{itemize}
            \item \textbf{多目标优化：}综合平衡后坐力、人机工效与价格
            \item \textbf{预算控制：}支持硬性上限或软性惩罚
            \item \textbf{商人等级过滤：}根据玩家等级筛选可购配件以及商人等级限制
            \item \textbf{自定义过滤器：}根据玩家需求，自定义过滤器，如必装槽位、必装配件、必装类别等，以及排除某些配件等
            \item \textbf{帕累托前沿：}忽略一个属性，可视化呈现其他属性之间的权衡曲线
            \item \textbf{枪匠任务：}自动求解枪匠任务，输出成本最优的改枪方案
        \end{itemize}
    \end{columns}
\end{frame}

% Section 2: Architecture
\section{系统架构}

% Slide: Pipeline Overview
\begin{frame}{整体流程}
    \begin{center}
    \begin{tikzpicture}[
        node distance=0.8cm,
        every node/.style={font=\small},
        box/.style={draw, rectangle, rounded corners, fill=blue!20, minimum width=2.4cm, minimum height=0.9cm, align=center},
        arrow/.style={->, thick, >=stealth, line width=1.5pt}
    ]
        % Main pipeline - horizontal
        \node[box] (fetch) {数据获取\\{\scriptsize Tarkov.dev API}};
        \node[box, right=of fetch] (compat) {构建兼容树\\{\scriptsize BFS遍历}};
        \node[box, right=of compat] (solve) {求解优化\\{\scriptsize CP-SAT}};
        \node[box, right=of solve, fill=green!25] (result) {最优方案};

        % User input from below
        \node[box, below=0.7cm of solve, fill=orange!25] (user) {用户约束};

        % Arrows
        \draw[arrow] (fetch) -- (compat);
        \draw[arrow] (compat) -- (solve);
        \draw[arrow] (solve) -- (result);
        \draw[arrow] (user) -- (solve);
    \end{tikzpicture}
    \end{center}
\end{frame}

% Slide 5: Fetching Data
\begin{frame}{数据获取}
    \begin{columns}
        \column{0.45\textwidth}
        \textbf{Tarkov.dev GraphQL API}

        \vspace{0.3cm}
        \begin{itemize}
            \item 获取全部武器、配件及兼容性规则
            \item \textbf{24小时内价格：}包含商人价格与跳蚤市场价格
            \item 商人等级要求与跳蚤市场解锁条件
            \item 本地缓存机制提升查询效率
        \end{itemize}

        \column{0.55\textwidth}
        \begin{center}
            \includegraphics[width=\textwidth]{imgs/tarkov_dev.jpg}
        \end{center}
    \end{columns}
\end{frame}

% Slide 6c: BFS Compatibility Tree
\begin{frame}{构建兼容性树}
    \begin{columns}
        \column{0.55\textwidth}
        \textbf{广度优先搜索（BFS）流程：}
        \begin{enumerate}
            \item 从基础武器的槽位出发
            \item 遍历每个槽位，获取所有可以安装在该槽位的配件
            \item 将这些配件加入待处理队列
            \item 继续探索每个配件\textit{自身}拥有的槽位
            \item 重复上述过程直至无新配件可以安装
        \end{enumerate}

        \column{0.45\textwidth}
        \begin{center}
        \begin{tikzpicture}[
            node distance=0.6cm,
            every node/.style={font=\scriptsize},
            slot/.style={draw, rectangle, rounded corners, fill=blue!15, minimum width=1.6cm, minimum height=0.4cm},
            item/.style={draw, rectangle, fill=green!15, minimum width=1.4cm, minimum height=0.4cm}
        ]
            \node[item] (gun) {M4A1};
            \node[slot, below left=0.5cm and -0.2cm of gun] (s1) {机匣槽};
            \node[slot, below right=0.5cm and -0.2cm of gun] (s2) {握把槽};
            \node[item, below=of s1] (upper) {上机匣};
            \node[slot, below=of upper] (s3) {枪管槽};
            \node[item, below=of s3] (barrel) {416mm枪管};

            \draw[->] (gun) -- (s1);
            \draw[->] (gun) -- (s2);
            \draw[->] (s1) -- (upper);
            \draw[->] (upper) -- (s3);
            \draw[->] (s3) -- (barrel);
        \end{tikzpicture}
        \end{center}
        \textit{\scriptsize 配件拥有槽位，槽位容纳配件}
    \end{columns}
\end{frame}

% Section 3: The Engine
\section{优化引擎}

% Slide 7: What is CP-SAT?
\begin{frame}{CP-SAT求解器简介}
    \textbf{CP-SAT} = 约束编程（Constraint Programming）+ 布尔可满足性（SAT）

    \vspace{0.3cm}
    \begin{block}{核心原理}
        Google OR-Tools库中的求解器，通过融合以下技术求解组合优化问题：
        \begin{itemize}
            \item \textbf{约束编程：}以声明式方式表达复杂规则
            \item \textbf{SAT求解：}高效的布尔逻辑推理引擎
            \item \textbf{线性规划：}优化数值目标函数
        \end{itemize}
    \end{block}

    \vspace{0.3cm}
    \textbf{为何适用于本问题？}
    \begin{itemize}
        \item 天然支持"选中/不选"的布尔决策建模
        \item 可处理数千个变量与约束条件
        \item 保证求得全局最优解，而非近似解
        \item 典型武器配置问题可在毫秒级完成求解
    \end{itemize}
\end{frame}

% % Slide 8: How CP-SAT Works
% \begin{frame}{CP-SAT工作原理}
%     \begin{columns}
%         \column{0.5\textwidth}
%         \textbf{1. 模型构建}
%         \begin{itemize}
%             \item 定义布尔变量：$X_i \in \{0, 1\}$
%             \item 添加约束条件
%             \item 设定目标函数
%         \end{itemize}

%         \vspace{0.3cm}
%         \textbf{2. 求解过程}
%         \begin{itemize}
%             \item 传播：推断确定值
%             \item 搜索：对未定变量进行分支
%             \item 回溯：遇到矛盾时返回
%             \item 剪枝：跳过不可行分支
%         \end{itemize}

%         \column{0.5\textwidth}
%         \textbf{3. 关键技术}
%         \begin{itemize}
%             \item \textit{单元传播：}约束确定的值立即应用
%             \item \textit{冲突学习：}记录失败原因以避免重复探索
%             \item \textit{界收紧：}利用目标值提前剪除次优分支
%         \end{itemize}

%         \vspace{0.3cm}
%         \textbf{效果：}高效遍历搜索空间，确保找到全局最优解。
%     \end{columns}
% \end{frame}

% Slide 9: Constraint Programming Overview
\begin{frame}{CP-SAT在武器配置中的应用}
    \textbf{核心思路：}明确定义合法配置的条件，由求解器自动搜索最优方案。

    \vspace{0.3cm}
    \begin{block}{模型三要素}
        \begin{enumerate}
            \item \textbf{决策变量：}为每个可用配件定义布尔变量$X_i$（$X_i = 1$表示选中）
            \item \textbf{约束条件：}兼容性规则、预算限制、必装槽位等
            \item \textbf{目标函数：}最大化人机工效、最小化后坐力或价格
        \end{enumerate}
    \end{block}

    \vspace{0.3cm}
    求解器遍历数百万种配件组合，实时剪除不合法配置，最终返回数学意义上的最优方案。
\end{frame}

% Slide 8: Constraint 1 - Mutex
\begin{frame}{约束1：槽位互斥}
    \textbf{规则：}每个槽位最多安装一个配件。

    \vspace{0.3cm}
    \begin{block}{数学表达}
        对于槽位$s$及其候选配件集合$\{i_1, i_2, ..., i_n\}$：
        \[
        \sum_{i \in \text{slot}_s} X_i \le 1
        \]
        其中$X_i = 1$表示配件$i$被选中，$X_i = 0$表示未选中。
    \end{block}

    \vspace{0.3cm}
    \textbf{示例：}M4的握把槽可安装GRAL-S\textit{或}Ergo PSG-1，但二者不可同时安装。
\end{frame}

% Slide 9: Constraint 2 - Dependency
\begin{frame}{约束2：父子依赖}
    \textbf{规则：}子配件的安装必须以父配件已安装为前提。

    \vspace{0.3cm}
    \begin{block}{数学表达}
        若配件$c$需安装于配件$p$所拥有的槽位：
        \[
        X_c \le X_p
        \]
        父配件未选中时，子配件不可选中。
    \end{block}

    \vspace{0.3cm}
    \textbf{示例：}安装前握把的前提是护木已安装在武器上，速瞄镜的前提是镜架已安装。

    \vspace{0.2cm}
    \begin{center}
    \textit{武器} $\to$ \textit{护木} $\to$ \textit{前握把}
    \end{center}
\end{frame}

% Slide 10: Constraint 3 - Conflicts
\begin{frame}{约束3：配件冲突}
    \textbf{规则：}部分配件存在互斥关系，不可同时安装。

    \vspace{0.3cm}
    \begin{block}{数学表达}
        对于API定义的冲突配件对$(a, b)$：
        \[
        X_a + X_b \le 1
        \]
        二者最多选其一。
    \end{block}

    \vspace{0.3cm}
    \textbf{示例：}AK以及M4的榴弹发射器和一些护木存在冲突，游戏不允许同时安装。
\end{frame}

% Slide 11: Constraint 4 - Required Slots
\begin{frame}{约束4：必装槽位}
    \textbf{规则：}特定槽位\textit{必须}安装配件，武器才可以正常使用。

    \vspace{0.3cm}
    \begin{block}{数学表达}
        对于必装槽位$s$：
        \[
        \sum_{i \in \text{slot}_s} X_i \ge 1
        \]
        该槽位必须安装至少一个兼容配件。
    \end{block}

    \vspace{0.3cm}
    \textbf{AR-15平台必装槽位：}
    \begin{itemize}
        \item 枪管、导气座、握把、拉机柄、机匣、护木
    \end{itemize}

    \vspace{0.2cm}
    缺少任一配件，武器均无法在游戏中正常使用。
\end{frame}

% Slide 12: Budget Constraint
\begin{frame}{约束5：预算限制}
    \textbf{规则：}配置总成本不得超过玩家设定的预算。

    \vspace{0.3cm}
    \begin{block}{数学表达}
        \[
        \sum_{i} \text{Price}_i \cdot X_i \le \text{Budget}
        \]
        所有选中配件的价格之和必须在预算范围内。
    \end{block}

    \vspace{0.3cm}
    \textbf{价格软约束机制（BFS阶段剪枝）：}
    \begin{itemize}
        \item 考虑商人等级（1-4级对应不同解锁价格）
        \item 对比跳蚤市场价格
        \item 自动选择每个配件的最低价购买渠道
    \end{itemize}
\end{frame}

% Slide 13: Objective Function
\begin{frame}{目标函数}
    \begin{block}{加权评分函数（最大化）}
        \[
        \text{Score} = W_E \cdot \text{人机} - W_R \cdot \text{后坐力} - W_P \cdot \text{价格}
        \]
        \vspace{-0.2cm}
        \begin{itemize}
            \item $W_E$：人机工效权重 \quad $W_R$：后坐力权重 \quad $W_P$：价格权重
        \end{itemize}
    \end{block}

    \vspace{0.2cm}
    \textbf{预设配置方案：}
    \begin{itemize}
        \item \textit{最高人机：}$W_E = 98\%, W_R = 1\%, W_P = 1\%$
        \item \textit{最低后坐力：}$W_E = 1\%, W_R = 98\%, W_P = 1\%$
        \item \textit{最低价格：}$W_E = 1\%, W_R = 1\%, W_P = 98\%$
        \item \textit{均衡方案：}$W_E = 34\%, W_R = 33\%, W_P = 33\%$
    \end{itemize}
\end{frame}

% Slide: UI - Optimizer
\begin{frame}{用户界面：优化配置}
    \begin{center}
        \includegraphics[width=0.9\textwidth]{imgs/page_optimize.png}
    \end{center}
\end{frame}

% Slide: Optimization Results - Recoil
\begin{frame}{优化结果：极限后坐力}
    \begin{center}
        \includegraphics[width=0.85\textwidth]{imgs/eg_recoil.png}
    \end{center}
\end{frame}

% Slide: Optimization Results - Ergo
\begin{frame}{优化结果：极限人机}
    \begin{center}
        \includegraphics[width=0.85\textwidth]{imgs/eg_ergo.png}
    \end{center}
\end{frame}

% Slide: Optimization Results - Balanced
\begin{frame}{优化结果：均衡方案}
    \begin{center}
        \includegraphics[width=0.85\textwidth]{imgs/eg_balanced.png}
    \end{center}
\end{frame}

% Slide 14: Stat Calculations
% \begin{frame}{属性计算规则}
%     \begin{columns}
%         \column{0.5\textwidth}
%         \textbf{人机工效}（加法叠加）：
%         \[
%         \text{人机}_{\text{最终}} = \text{人机}_{\text{基础}} + \sum_{i} \text{人机}_i
%         \]
%         各配件的人机加成直接累加。

%         \vspace{0.5cm}
%         \textit{游戏内存在100点软上限。}

%         \column{0.5\textwidth}
%         \textbf{后坐力}（乘法叠加）：
%         \[
%         \text{后坐力}_{\text{最终}} = \text{后坐力}_{\text{基础}} \times \left(1 + \sum_{i} r_i\right)
%         \]
%         其中$r_i$为百分比修正值（如$-0.05$表示$-5\%$）。

%         \vspace{0.3cm}
%         \textit{多个减后坐力配件叠加存在边际递减效应。}
%     \end{columns}
% \end{frame}

% Section 4: Advanced Scalarization: Sweet Spot Mode
\section{高级标量化：甜点模式 (Sweet Spot)}

% Slide: Limitation of Linear Weighted Sum
\begin{frame}{传统线性加权求和的局限性}
    \textbf{传统目标函数：} $\max \; W_E \cdot \text{人机} - W_R \cdot \text{后坐力} - W_P \cdot \text{价格}$
    
    \vspace{0.3cm}
    \begin{columns}
        \column{0.52\textwidth}
        \textbf{非凸帕累托前沿问题：}
        \begin{itemize}
            \item 线性加权求和在目标空间中以\textbf{直线超平面}进行扫描。
            \item 塔科夫配件改装空间具有离散组合特性，存在明显的\textbf{非凸凹陷区域}（如重型消音器 vs 轻量制退器的跳跃权衡）。
            \item \textbf{数学盲区：}位于非凸凹陷内部的配置点，在\textbf{任意线性权重组合} $(W_E, W_R, W_P)$ 下都无法被求解器触达！
        \end{itemize}

        \column{0.48\textwidth}
        \begin{block}{实证分析数据}
            对常见基准武器（M4A1、AK-74N、MP7A2）进行穷举分析发现：
            \begin{itemize}
                \item 真实帕累托前沿中 \textbf{45\%--61\%} 的非支配最优配置位于非凸凹陷区！
                \item 线性权重会直接从极限后坐力跳跃到极限人机，跳过最实用的\textbf{“甜点改装”}。
            \end{itemize}
        \end{block}
    \end{columns}
\end{frame}

% Slide: The 3D Ideal Reference Point
\begin{frame}{3D 理想参考点 (Ideal Point)}
    \textbf{核心思路：}放弃任意斜率的线性权重，转而度量与\textbf{乌托邦理想目标点}的差距。
    
    \vspace{0.3cm}
    \begin{columns}
        \column{0.5\textwidth}
        \textbf{1. 乌托邦理想点 ($\mathbf{z}^*$)：}
        \begin{itemize}
            \item $z^*_E = \max_{\mathbf{x}} \text{Ergo}(\mathbf{x})$ \quad (单目标最高人机)
            \item $z^*_R = \min_{\mathbf{x}} \text{Recoil}(\mathbf{x})$ \quad (单目标最低后坐力)
            \item $z^*_P = \min_{\mathbf{x}} \text{Price}(\mathbf{x})$ \quad (单目标最低价格)
        \end{itemize}
        \textit{通过3次独立的单目标求解在毫秒内获得。}

        \column{0.5\textwidth}
        \textbf{2. 归一化目标差距 ($g_i$)：}
        对各个目标维度 $i \in \{E, R, P\}$：
        \[
        g_i(\mathbf{x}) = \frac{|f_i(\mathbf{x}) - z^*_i|}{\text{Range}_i} \in [0, 1]
        \]
        \begin{itemize}
            \item $g_i = 0$：达到理论最佳极限。
            \item 消除人机点数、后坐力百分比与卢布价格之间的量纲差异。
        \end{itemize}
    \end{columns}
\end{frame}

% Slide: Augmented Tchebycheff Formulation
\begin{frame}{增广切比雪夫标量化 (Augmented Tchebycheff)}
    \textbf{数学建模（甜点模式）：}
    \[
    \min_{\mathbf{x}, \alpha} \left( \alpha + \rho \sum_{i \in \{E, R, P\}} \lambda_i \cdot g_i(\mathbf{x}) \right)
    \]
    \[
    \text{约束条件：} \quad \lambda_i \cdot g_i(\mathbf{x}) \le \alpha, \quad \forall i \in \{E, R, P\}
    \]

    \vspace{0.2cm}
    \begin{columns}
        \column{0.5\textwidth}
        \textbf{算法优势：}
        \begin{itemize}
            \item $\alpha = \max_i (\lambda_i g_i(\mathbf{x}))$ 最小化\textbf{最大维度的短板}，促进水桶级均衡。
            \item 增广项 $\rho \sum \lambda_i g_i$ ($\rho = 10^{-4}$) 严格保证\textbf{强帕累托最优}，剔除平坦弱解。
        \end{itemize}

        \column{0.5\textwidth}
        \textbf{工程实现特性：}
        \begin{itemize}
            \item 仅引入 \textbf{1 个连续辅助变量} ($\alpha$)。
            \item \textbf{不增加任何 0-1 整数变量}，求解时间保持在 $< 15\text{ms}$。
            \item 100\% 挖掘非凸帕累托甜点改装方案。
        \end{itemize}
    \end{columns}
\end{frame}

% Slide: Geometric Intuition: Hyperplanes vs L-Contours
\begin{frame}{几何直观：超平面 vs. $L_\infty$ 等值线}
    \begin{center}
    \begin{tikzpicture}[scale=0.82]
        % Axes
        \draw[->, thick] (0,0) -- (6.5,0) node[right] {\small 后坐力/价格差距 ($g_R$)};
        \draw[->, thick] (0,0) -- (0,4.6) node[above] {\small 人机差距 ($g_E$)};

        % Non-convex Pareto front
        \draw[very thick, blue] (1.0, 4.0) to[out=-50, in=165] (2.6, 2.3) to[out=-15, in=125] (4.2, 2.0) to[out=-55, in=170] (5.8, 0.7);

        % Ideal Point z*
        \node[fill=green!70!black, circle, inner sep=2.5pt, label={[font=\footnotesize, green!70!black]below left:$\mathbf{z}^* = (0, 0)$ (理想点)}] at (0.2, 0.2) {};

        % Linear Hyperplane tangent touching extremes, missing the knee
        \draw[dashed, red, thick] (0.2, 4.8) -- (6.6, -0.2) node[right, font=\scriptsize, red] {线性超平面};
        \node[fill=red, circle, inner sep=2pt, label={[font=\scriptsize]above right:极端配置 A}] at (1.0, 4.0) {};
        \node[fill=red, circle, inner sep=2pt, label={[font=\scriptsize]above right:极端配置 B}] at (5.8, 0.7) {};

        % Tchebycheff L-contour
        \draw[thick, purple] (2.6, 4.2) -- (2.6, 2.3) -- (6.0, 2.3) node[above, font=\scriptsize, purple] {$L_\infty$ 等值线};
        \node[fill=purple, star, star points=5, inner sep=2.5pt, label={[font=\scriptsize, purple]above right:\textbf{甜点配置！}}] at (2.6, 2.3) {};

        \node[draw, rounded corners, fill=yellow!15, font=\scriptsize, align=center] at (3.8, 3.4) {非凸凹陷区\\(线性加权无法触达)};
    \end{tikzpicture}
    \end{center}
    \textbf{几何结论：}直线切平面只能触及外凸极值点，而切比雪夫 L 型角等值线能深入非凸凹陷中，精准定位均衡甜点解。
\end{frame}

% Section 5: Aiming Physics & EvoErgo Optimization
\section{真实开镜物理与 EvoErgo 优化}

% Slide: In-Game ADS Physics Problem
\begin{frame}{游戏内开镜机制与过冲甩头 (Overswing)}
    \textbf{实战中的物理矛盾：}
    \begin{itemize}
        \item 高人机工效可加快开镜瞄准（ADS）过渡速度并降低体力消耗。
        \item 但枪械重量过大会带来\textbf{转动惯性与晃动}。
    \end{itemize}

    \vspace{0.3cm}
    \begin{columns}
        \column{0.5\textwidth}
        \begin{block}{什么是瞄准过冲 (Overswing)？}
            当枪械重量超过当前人机工效的承载阈值时，快速开镜会\textbf{甩过准星中心点并反向晃动}。
            \begin{itemize}
                \item 导致第一发定位延迟，极度影响近战胜率。
                \item 100 人机、7.5\,kg 的全装武器会出现剧烈过冲晃动。
                \item 65 人机、3.2\,kg 的轻量化武器能做到瞬间平稳锁头。
            \end{itemize}
        \end{block}

        \column{0.5\textwidth}
        \begin{center}
        \textit{"脱离重量单独看人机数值是失真的。"}
        
        \vspace{0.3cm}
        \textbf{社区实测研究：}\\
        SpaceMonkey37 与 EFTForge 建立了人机、重量与过冲的实证非线性模型。
        \end{center}
    \end{columns}
\end{frame}

% Slide: Quadratic EvoErgoDelta (EED)
\begin{frame}{二次非线性 EvoErgoDelta ($\text{EED}$) 指标}
    \textbf{社区公式 (SpaceMonkey37 / EFTForge)：}
    \[
    \text{EED} = \left[ 100 - 0.015 \cdot (100 - \text{人机})^2 \right] - 15 \cdot \text{重量}_{\text{kg}}
    \]

    \vspace{0.2cm}
    \begin{columns}
        \column{0.5\textwidth}
        \textbf{$\text{EED}$ 的物理意义：}
        \begin{itemize}
            \item \textcolor{green!60!black}{\textbf{$\text{EED} \ge 0$（无开镜过冲）}}：重量在人机承载阈值以内，快速开镜指哪打哪。
            \item \textcolor{red!80!black}{\textbf{$\text{EED} < 0$（存在开镜过冲）}}：转动惯性过大，开镜后产生枪口晃动。
        \end{itemize}

        \column{0.5\textwidth}
        \textbf{线性近似分值 (EvoErgo)：}
        \[
        \text{EvoErgo} = \min(100, \text{人机}) - 15 \cdot \text{重量}_{\text{kg}}
        \]
        \begin{itemize}
            \item 结果面板展示的标准社区分值。
            \item 以 1:15 kg 的交换比直观衡量人机与重量的综合效率。
        \end{itemize}
    \end{columns}
\end{frame}

% Slide: Linearizing Quadratic EED via Tangent k-Sweep
\begin{frame}{二次 EED 的线性化求解：切线 $k$-Sweep 算法}
    \textbf{求解难点：}$\text{EED}$ 包含二次非线性项，但求解器需要线性目标函数以实现毫秒级响应。

    \vspace{0.3cm}
    \textbf{切线交换率：}边际人机价值随当前人机点数 $E$ 动态变化：
    \[
    k(E) = \frac{\partial \text{EED}}{\partial E} = 0.03 \cdot (100 - E) + 15 \quad (\text{切线斜率 } \Delta W / \Delta E)
    \]

    \vspace{0.2cm}
    \begin{block}{切线 $k$-Sweep 候选求解算法 (\texttt{solveEvoErgo})}
        \begin{enumerate}
            \item 在基准交换率 $k \in \{15.0, 10.0, 7.5, 5.6\}$ 及 $k(E_0)$ 上并行运行线性求解。
            \item 获得多个满足整数约束的候选配置 $\mathbf{x}^{(1)}, \mathbf{x}^{(2)}, \dots, \mathbf{x}^{(m)}$。
            \item 对所有候选解计算精确的二次 $\text{EED}(\mathbf{x}^{(j)})$ 指标。
            \item 在 $< 10\text{ms}$ 内返回使二次 $\text{EED}$ 全局最大的改装方案。
        \end{enumerate}
    \end{block}
\end{frame}

% Slide: Prevent Overswing Hard Constraint
\begin{frame}{硬性约束：防止开镜过冲 (Prevent Overswing)}
    \textbf{硬性要求：}确保武器在任何情况下绝不过冲（$\text{EED} \ge 0$）。
    \[
    \text{重量}_{\text{kg}} \le \text{阈值}(E_{\text{eff}}) = \frac{100 - 0.015 \cdot (100 - E_{\text{eff}})^2}{15}
    \]

    \vspace{-0.1cm}
    \begin{columns}
        \column{0.55\textwidth}
        \begin{center}
        \begin{tikzpicture}[scale=0.75]
            \draw[->, thick] (0,0) -- (5.5,0) node[right] {\scriptsize $E_{\text{eff}}$};
            \draw[->, thick] (0,0) -- (0,4.2) node[above] {\scriptsize $W$ (kg)};

            \fill[green!15, opacity=0.8] (0.5, 0) -- (0.5, 0.6) to[out=40, in=195] (3.0, 2.3) to[out=15, in=180] (5.0, 3.3) -- (5.0, 0) -- cycle;
            \draw[very thick, teal] (0.5, 0.6) to[out=40, in=195] (3.0, 2.3) to[out=15, in=180] (5.0, 3.3);

            \node[font=\tiny, green!50!black, align=center] at (2.8, 1.0) {\textbf{零过冲安全区}\\($\text{EED} \ge 0$)};
            \node[font=\tiny, red!70!black, align=center] at (1.8, 3.2) {\textbf{过冲晃动区}\\($\text{EED} < 0$)};

            \draw[dashed, thick, orange] (1.6, 1.4) -- (4.4, 3.2) node[above, font=\tiny, orange] {切线割平面};
            \node[fill=blue, circle, inner sep=1.5pt] at (3.0, 2.3) {};
        \end{tikzpicture}
        \end{center}

        \column{0.45\textwidth}
        \textbf{切线割平面线性化：}
        \begin{itemize}
            \item 用一阶切线不等式包络非线性曲线：
            \[
            W \le T(E_0) + T'(E_0)(E - E_0)
            \]
            \item 凸上界剪枝，无需引入额外 0-1 变量，求解时间 $< 15\text{ms}$。
        \end{itemize}
    \end{columns}
\end{frame}

% Slide: Equipment Ergonomics Modifier (b)
\begin{frame}{装备人机工效惩罚系数 ($b$)}
    \textbf{实战环境考量：}防弹衣、防弹弹挂、战术头盔和大型背包会大幅削减有效人机工效。

    \vspace{0.3cm}
    \begin{block}{有效人机工效公式}
        \[
        E_{\text{eff}} = E_{\text{weapon}} \cdot (1 + b), \quad b \in [-0.40, \; 0.00]
        \]
        \begin{itemize}
            \item $b$：防具与装备的总人机惩罚百分比。
            \item 例：重型6级甲 Zabralo ($b = -28\%$) + 阿尔金头盔 ($b = -7\%$) $\implies b = -35\%$。
        \end{itemize}
    \end{block}

    \vspace{0.3cm}
    \textbf{实战价值：}求解器自动计算装备削减后的有效人机并动态收紧过冲重量阈值，保证全装实战中开镜指哪打哪。
\end{frame}

% Section 6: Key Features
\section{特色功能}
% Slide: Gunsmith Task Solver
\begin{frame}{枪匠任务求解器}
    \textbf{背景：}枪匠任务要求玩家按照指定规格改装武器，对新手玩家具有一定难度。

    \vspace{0.2cm}
    \begin{columns}
        \column{0.5\textwidth}
        \textbf{任务定义（JSON格式）：}
        \begin{itemize}
            \item 指定武器型号
            \item 属性要求：
            \begin{itemize}
                \item 人机工效下限、后坐力上限
                \item 弹匣容量、重量、瞄准距离等
            \end{itemize}
            \item 必须安装的特定配件
            \item 必须包含的配件类别（如消音器）
        \end{itemize}

        \column{0.5\textwidth}
        \textbf{求解流程：}
        \begin{enumerate}
            \item 加载任务配置文件
            \item 将任务要求转化为约束：
            \item 目标函数：最小化总价格
            \item 输出成本最低的达标方案
        \end{enumerate}
    \end{columns}

    \vspace{0.2cm}
    \textbf{效果：}可自动求解全部28个枪匠任务，输出成本最优的通关方案。
\end{frame}

% Slide: Gunsmith List
\begin{frame}{枪匠任务列表}
    \begin{center}
        \includegraphics[width=0.9\textwidth,height=0.75\textheight,keepaspectratio]{imgs/eg_gunsmith.png}
    \end{center}
\end{frame}

% Slide: Gunsmith Details
\begin{frame}{枪匠任务求解详情}
    \begin{center}
        \includegraphics[width=0.9\textwidth,height=0.75\textheight,keepaspectratio]{imgs/eg_gunsmith_2.png}
    \end{center}
\end{frame}

% Slide 9: Pareto Frontier
\begin{frame}{帕累托前沿分析}
    \textbf{概念说明：}
    \begin{itemize}
        \item 帕累托前沿展示了两个冲突目标（如后坐力与人机工效）之间的最优权衡曲线。
    \end{itemize}

    \begin{center}
    \begin{tikzpicture}[scale=0.85]
        \draw[->] (0,0) -- (6,0) node[right] {\small 后坐力};
        \draw[->] (0,0) -- (0,4) node[above] {\small 人机};

        \draw[thick, blue] (0.5, 1) to[out=45, in=180] (2.5, 2.8) to[out=0, in=135] (5, 3.3);

        \node[fill=red, circle, inner sep=2pt, label={[font=\scriptsize]left:低后座}] at (0.5, 1) {};
        \node[fill=green, circle, inner sep=2pt, label={[font=\scriptsize]below right:均衡}] at (2.5, 2.8) {};
        \node[fill=red, circle, inner sep=2pt, label={[font=\scriptsize]right:高人机}] at (5, 3.3) {};

        \node[font=\small] at (1.2, 3.5) {不可达};
        \node[font=\small] at (4.5, 1) {非最优};
    \end{tikzpicture}
    \end{center}

    \textbf{应用价值}：帮助玩家直观判断后坐力与人机工效之间的权衡，选择最适合自己的配置方案。
\end{frame}

% Slide: Pareto Visualization
\begin{frame}{帕累托前沿可视化}
    \begin{center}
        \includegraphics[width=0.9\textwidth,height=0.75\textheight,keepaspectratio]{imgs/page_pareto.png}
    \end{center}
\end{frame}

% Slide: Future Work
\begin{frame}{尚未实现的功能}
    \begin{itemize}
        \item \textbf{武器隐藏属性}：目前无法获取武器的隐藏属性，如后坐力角度等
        \item \textbf{商人以物易物}：目前无法获取商人以物易物的物品列表，以及商人以物易物的价格
    \end{itemize}
\end{frame}

% Section 7: Conclusion
\section{总结}
% Acknowledgements
\begin{frame}{致谢}
    \begin{center}
        \textbf{\Large 数据来源}

        \vspace{0.2cm}
        \includegraphics[height=0.9cm]{logos/tarkov-dev-logo.png}

        \vspace{0.1cm}
        社区开源GraphQL API

        \vspace{0.4cm}
        \textbf{\Large AI编程助手}

        \vspace{0.3cm}
        \begin{tabular}{c @{\hspace{1cm}} c @{\hspace{1cm}} c}
            \includegraphics[height=1.1cm]{logos/claude-color.png} &
            \includegraphics[height=1.1cm]{logos/gemini-color.png} &
            \includegraphics[height=1.1cm]{logos/minimax-color.png} \\[0.2cm]
            \textbf{Claude Code} & \textbf{Gemini CLI} & \textbf{MiniMax} \\[0.1cm]
            Anthropic & Google & MiniMax \\
        \end{tabular}

        \vspace{0.3cm}
        \textit{感谢以上AI助手在代码编写、调试及文档撰写方面提供的支持}
    \end{center}
\end{frame}

\end{document}
